% ============================================================
\section{Solver Layer: Exactness, Green Operators, and Homotopy}
\label{subsec:solver-layer}
% ============================================================

The preceding chapters require only projection algebra and additive closure.
Sometimes, however, the observable operator data sits inside a finite
Hilbert-complex package. In that case an additional solver layer becomes
available: exactness, Green operators, and canonical homotopies.

\begin{definition}[Finite Hilbert complex]
\label{def:hilbert-complex}\lean{finiteHilbertComplex}
A \emph{finite Hilbert complex} is a sequence of finite-dimensional inner
product spaces and linear maps
\[
0 \longrightarrow H^0 \xrightarrow{D_0} H^1 \xrightarrow{D_1} \cdots
\xrightarrow{D_{m-1}} H^m \longrightarrow 0
\]
such that
\[
D_{k+1} D_k = 0
\qquad
\text{for every } k.
\]
\end{definition}

\begin{definition}[Exactness]
\label{def:exactness}\lean{complexExactness}
The complex is \emph{exact at degree $k$} if
\[
\Ker(D_k) = \Img(D_{k-1}).
\]
Equivalently, the $k$th cohomology group vanishes:
\[
\CCCoh{k}(D) = 0.
\]
\end{definition}

\begin{definition}[Hodge Laplacian and Green operator]
\label{def:green-operator}\lean{hodgeLaplacian, greenOperator}
At degree $k$, define the Hodge Laplacian by
\[
\CCLap{k} := D_{k-1} D_{k-1}^* + D_k^* D_k.
\]
When $\CCLap{k}$ is invertible, its inverse
\[
\CCGreen{k} := \CCLap{k}^{-1}
\]
is the \emph{Green operator} at degree $k$.
\end{definition}

\begin{lemma}[Intertwining of the Laplacian with the differential]
\label{lem:laplacian-intertwining}\lean{laplacian_intertwines_differential}
For every degree $k$,
\[
D_k \CCLap{k} = \CCLap{k+1} D_k.
\]
Consequently, whenever the Green operators exist,
\[
D_k \CCGreen{k} = \CCGreen{k+1} D_k.
\]
\end{lemma}

(Proof in Appendix~\ref{sec:appendix-elimination}.)

\begin{theorem}[Homotopy identity with adjacent Green operators]
\label{thm:homotopy-identity}\lean{homotopyIdentity}
Assume the Green operators \(\CCGreen{k}\) and \(\CCGreen{k+1}\) both exist.
Define
\[
\CCHom{k} := D_{k-1}^* \CCGreen{k},
\qquad
\CCHom{k+1} := D_k^* \CCGreen{k+1}.
\]
Then
\[
D_{k-1} \CCHom{k} + \CCHom{k+1} D_k = I_{H^k}.
\]
\end{theorem}

(Proof in Appendix~\ref{sec:appendix-elimination}.)

\begin{remark}[Lean boundary for the solver layer]
\label{rem:solver-layer-lean-boundary}
\lean{homotopyIdentity}
On the active Lean surface, the homotopy identity is treated as an explicit
solver-interface property rather than as a bedrock theorem forced from
exactness alone. The theorem above is therefore one sufficient realization of
that interface when adjacent Green operators are available, not part of the
primitive calculus itself.
\end{remark}

\begin{definition}[Holographic saturation]
\label{def:holographic-saturation}\lean{holographicSaturation}
Let $R_k : H^k \to B^k$ be a boundary restriction map. The degree-$k$ layer is
\emph{holographically saturated} if
\[
\Ker(R_k) = \Img(D_{k-1}).
\]
Equivalently, the relative cohomology at degree $k$ vanishes:
\[
\CCCohRel{k}{D,R} = 0.
\]
\end{definition}

\paragraph{Scope of the solver layer.}
This layer is optional. It is not part of the bedrock axioms of the calculus.
Rather, it is the exact package that becomes available once the operator data
is known to form an exact finite Hilbert complex.
